% ============================================================
%  Harald Høffding, "Religion og Videnskab" — "Religion and Science"
%  Religionshistoriske Smaaskrifter I
%  Copenhagen: Gyldendal — Nordisk Forlag, 1910
%  TRANSLATION (English)
%  Mirrors transcription.tex 1:1. Source of truth = transcription.tex.
%  Progress: printed pp. 1–5 (PDF 10–14) translated; continues mid-work.
% ============================================================
\documentclass[12pt,a4paper]{article}
\usepackage[utf8]{inputenc}
\usepackage[T1]{fontenc}
\usepackage{libertinus}
\usepackage{libertinust1math}
\usepackage{textalpha}
\usepackage[english]{babel}
\usepackage[protrusion=true,expansion=true]{microtype}
\usepackage{setspace}
\usepackage[margin=1.2in]{geometry}
\usepackage{fancyhdr}
\usepackage{enumitem}
\usepackage{hyperref}

\hypersetup{
  pdftitle={Religion and Science},
  pdfauthor={Harald Høffding},
  hidelinks
}

\pagestyle{fancy}
\fancyhf{}
\rhead{Høffding, \emph{Religion and Science} (1910)}
\cfoot{\thepage}

\setstretch{1.3}

\title{\textbf{Religion and Science}\\[0.4em]
\large \emph{Religionshistoriske Smaaskrifter} I}
\author{Harald Høffding\\[4pt]
  \small Translated from the Danish}
\date{Copenhagen: Gyldendal --- Nordisk Forlag, 1910}

\begin{document}
\maketitle
\thispagestyle{fancy}

\bigskip

\noindent Within every culture we know, a critical condition sets in once
the division of labour comes into force. In the material sphere it then
becomes no longer possible for the individual worker to produce a work in
which he can have the feeling of having worked with all his powers. What
the individual can do now becomes, most often, only the production of a
single part of a whole, a part that requires merely the application of a
single faculty. The development of his power and his talent thereby becomes
one-sided. Like his work, his very personality readily becomes but a
fragment. In the spiritual sphere the division of labour means that art,
science, religion, and morality, which in earlier times formed a unity, now
step apart and often step into opposition to one another as contending
powers. In earlier cultural periods religion stood as the great central
current of life, to which all spiritual forces and drives gave their
contribution, and which on the other hand yielded full satisfaction for the
artistic and the intellectual need, in so far as such need stirred at all.
But gradually the spiritual needs seek their satisfaction along different
paths. The investigative urge is no longer satisfied by the answers that
religious ideas can give to the manifold questions that life and experience
may raise. Special methods are developed, special forms for posing the
questions and for seeking the answers. The artistic impulse likewise goes
more and more its own ways, and an independent foundation is sought for
moral valuation. History shows that, in contrast to these new formations,
religion has a tendency to hold fast to the notions and modes of action
once established. It resists the spiritual division of labour as long as it
can, and it seeks as long as possible to close its eyes to the altered
conditions of spiritual life. From this follows a dividedness and a
splintering in spiritual life. There is lacking that gathering and harmony,
that concentration of all the forces of the mind in one common direction,
which under the earlier conditions was comparatively easy to attain. On
this rests the fact that there exists a religious problem, and that such a
problem will continue to exist until the spiritual life of individuals and
of the race as a whole has won new gathered and gathering forms for the
satisfaction of its innermost need---if indeed such new forms are possible
at all.

I shall here bring forward only a few features of the intellectual side of
the problem. The whole question is, for me, not exhausted by being viewed
from this side. But it is of great importance to become clear about what
the opposition that again and again asserts itself between religion and
science really rests upon. Often one finds---and this holds of both parties
that here stand over against one another---the ground of this opposition
put in the wrong place. There is much superficial and dogmatic harping on
science's real or supposed results on the one side, and there is on the
other side much short-sightedness with respect to the change that the life
of thought has in fact undergone, quite apart from the many conspicuous
results. In the collisions to which the opposition leads, the onlooker does
not miss honesty in the confessions made from the one side or the other.
For nowadays there is, after all, comparatively little danger in making
one's confession, in whatever direction it may go. By contrast, one has
ample occasion to miss that intellectual integrity which does not let the
mind rest before the great questions have been examined from all sides and
by all the means that stand at one's disposal. This is unfortunately not a
virtue that religious ethics has especially inculcated, and many of
religion's opponents show, by the character of their polemic, that they too
do not know the strict demands it makes. The want that an onlooker may feel
here is bound up with the fact that religious debate in our days has mostly
the character of agitation. On both sides it is a matter of winning
adherents, of dissemination, not of deepening. New thoughts, accordingly,
this modern debate has not brought either.

What I wish to bring here is likewise not new thoughts, but old thoughts
that have not been allowed to come to their full due.

\begin{center}---\end{center}

If we go back to times in which religion was not merely help and
consolation, but also gave answers to what thought could then ask
about---times, that is, in which the whole conception of the world had a
religious character---then there are two peculiarities that especially come
to the fore. The whole conception of the world has an intuitive character.
Existence stands as a whole whose several regions can be seen or at least
represented, without its being necessary to break with what the senses
immediately seemed to teach. Cosmic space, as it presented itself to
sensation and imagination, contained the possibility of localizing all the
great events with which religious proclamation dealt. And when
changes---especially momentous changes---entered into this world of
intuitability, it was not by the arduous path of research that \emph{the
explanation} was to be sought. Such changes were explained by the
intervention of personal beings in the course of events. It is
intuitability, and it is the principle of personal causes (what in the
widest sense may be called animism), that characterizes a specifically
religious conception of the world.

But both the concept of space and the concept of cause have, under the
influence of the independent urge to knowledge, undergone changes that are
here of decisive importance.

\medskip

1. We have a natural urge to refer everything that happens to definite
places in space, and religion in its classical periods does not deny this
urge. Religious representations have the character of intuition, not of
concept, and their great influence, their power to grip and to fill, rests
in no small part upon this. The content of the higher popular religions is
a series of historical events into which the fate of men and of human life
is woven---a great drama, and a drama played out upon a stage,
sense-intuition, which imagination can in any case very well encompass
without bursting. The whole of existence, so far as it is accessible at
this stage, is taken up into this stage, and its various regions furnish
places for the various scenes in the drama. The earth stands as centre, or
as the lowest place, and from the visible heaven, which is the world of the
deity, divine beings descend to intervene in the fate of men, perhaps
themselves to suffer and die, before they ascend again into the heavenly
regions, to which they thereupon open access for those who will follow
them. ``All men,'' says Aristotle, ``suppose that there are gods, and all,
both Hellenes and barbarians, assign to the deity the highest place.'' This
utterance applies to all popular religions, however different they may
otherwise be. Religious representations were woven in a natural way into
men's whole involuntarily formed conception, and there as yet existed no
science that necessitated a radical change of that conception. There was no
vagueness or indistinctness in the representation of the divine; this could
be localized like everything else; one could point toward the place from
which the divine forces welled forth, as well as to the place at which they
worked. The divine thoughts were, to be sure, raised as high above man's
thoughts as heaven above earth; but heaven was precisely not raised so high
above earth that intuition had to give up. Despite the difference in power
and dignity, a vivid and intimate relation to the gods was nevertheless
possible.
% [translation continues: printed p. 6 / PDF p. 15 — "Ingen Folkereligion har forladt dette Verdensbillede…"]

\end{document}
